\pagestyle{empty}
%
% Deckblatt
%
\begin{figure}[p]
\begin{tikzpicture}
	% Rahmen
	\draw[rounded corners=4mm] (151mm,213mm) -- (151mm,196mm) -- (0mm,196mm) -- (0mm,145mm) -- (151mm,145mm) -- (151mm,63mm);
	% F
	\draw (62mm,3mm) -- (62mm,59mm) -- (102mm,59mm) (62mm,31mm) -- (83mm,31mm);
	% A
	\draw (73mm,3mm) -- (106mm,58mm) -- (106mm,3mm) (90mm,31mm) -- (106mm,31mm);
	% U
	\draw (113mm,58mm) -- (113mm,18mm) arc (180:360:19mm)  -- (151mm,58mm);

	% FAU Logo
	\node at (25mm,175mm){\includegraphics[keepaspectratio=true,width=35mm]{deckblatt/FAU-Siegel.png}};
	% Arbeitsart
	\node at (25mm,151mm){\Large\thethesis};
	% Titel
	\node at (97mm,175mm){
		\parbox{10cm}{
			\centering \bfseries \itshape \Large \thesisTitleFormatted
		}
	};
	% Autor
	\node at (97mm,151mm){
			\bfseries \itshape \Large \myname
	};
	
	\node at (60mm,59mm)[below left,inner sep=0]{
		\parbox{6cm}{
			\begin{flushright}
				\large Lehrstuhl für Informatik 6 \\ (Datenmanagement)
			\end{flushright}
		}
	};
	\node at (60mm,25mm)[above left]{
		\parbox{6cm}{
			\begin{flushright}
				\large Department Informatik
			\end{flushright}
		}
	};
	\node at (60mm,-2mm)[above left,inner sep=0]{
		\parbox{6cm}{
			\begin{flushright}
				\large Friedrich Alexander-\\Universität\\Erlangen-Nürnberg
			\end{flushright}
		}
	};
\end{tikzpicture}
\end{figure}
%
% Titelseite
%
\begin{titlepage}

  \begin{center}
    
    {\Huge \bfseries
      \thesistitle{}\\
    } 
    
    \vspace*{1cm}
    \thethesis{} im Fach Informatik
    \vspace{2cm}
    
    {\large vorgelegt von} \\
    \vspace*{0.7cm}
    {\Large \bfseries \myname} \\
    \vspace*{0.7cm}
    {\large geb. \birthday{} in \birthplace{}} 
    
    \vspace{1.5cm}
    
    angefertigt am 

    \vspace{1cm}
    
    {\bfseries
      Lehrstuhl für Informatik 6 (Datenmanagement) \\
      Department Informatik \\
      Friedrich-Alexander-Universität Erlangen-Nürnberg (FAU)\\
      }
    
    \vspace{0.5cm}
\end{center}
\begin{tabbing}
    Betreuer: \= \corrector{}\\
    \> \tutor{} 
\end{tabbing}
    \vspace{0.25cm}
    
\begin{tabbing}
  Beginn der Arbeit: \startofwork{} \\
  Abgabe der Arbeit: \dayofdoom{}
\end{tabbing}

\end{titlepage}

\begin{comment}
%
% Leere Seite für doppelseitiges Layout einfügen
\cleardoubleemptypage
%
% Leeres Seitenlayout: Keine Seitennummern, kein Header oder Footer
%
Ich versichere, dass ich die Arbeit ohne fremde Hilfe und ohne Benutzung
anderer als der angegebenen Quellen angefertigt habe und dass diese Arbeit in
gleicher oder ähnlicher Form noch keiner anderen Prüfungsbehörde
vorgelegen hat und von dieser als Teil einer Prüfungsleistung angenommen
wurde. Alle Ausführungen, die wörtlich oder sinngemäß übernommen
wurden, sind als solche gekennzeichnet.

\vspace{2cm}

\noindent
Der Friedrich-Alexander-Universität Erlangen-Nürnberg (FAU), vertreten durch
den Lehrstuhl für Informatik 6 (Datenmanagement), wird für Zwecke der
Forschung und Lehre ein einfaches, kostenloses, zeitlich und örtlich
unbeschränktes Nutzungsrecht an den Arbeitsergebnissen des Masterprojektes
einschließlich etwaiger Schutzrechte und Urheberrechte eingeräumt.

\vspace{2cm}
Erlangen, den \dayofdoom{}

\vspace{2cm}
\myname{} \hfill \ 

\end{comment}

\vspace{0,5cm}
%
% Leere Seite für doppelseitiges empty Layout einfügen
\cleardoubleemptypage
%
% Seitenlayout mit Seitennummern, allerdings kein Header oder Footer
\pagestyle{plain}
%
% Römische Seitenzahlen
\pagenumbering{roman}
% Da wir 6 Seiten ohne Seitennummerierung im Seitenlayout empty verwenden, reicht ein 
% \usepackage[pdftex,pdfpagelabels,plainpages=false,]{hyperref} nicht mehr aus. Es muss zusätzlich noch der 
% Seitenzähler auf 7 anstatt wie defaultmäßig auf 1 für die folgend Seiten gesetzt werden.
\setcounter{page}{7}
%
\chapter*{Abstract}
\section*{\thesistitleGB}
\selectlanguage{english}

The Open Source Learning Management System ILIAS is a an extensive tool to support students at universities learning. One big feature of ILIAS is the testing module. Using this lecturers are enabled to give their students a test they have to solve before accessing important documents. Of course the testing module can be used to create real (online) exams, too.

While the testing module supports a bunch of different question types by default, including Single Choice, Multiple Choice and Ordering questions, more complex question types have to be added by installing a plugin.

In this project - called \eigenname{\textbf{Q}uestion \textbf{P}lugin for \textbf{I}lias - \textbf{SQL} \textbf{(QPI-SQL)}} - we developed an ILIAS plugin, named \eigenname{assSQLQuestion}. The plugin enables lecturers to add interactive SQL questions to their tests. Students can execute their answers and their work is graded automatically by applying different metrics.
 
\clearpage{\pagestyle{empty}\cleardoublepage}
\chapter*{Kurzfassung}
\section*{\thesistitle}
\selectlanguage{ngerman}

Das Open Source Lernmanagement System ILIAS ist ein umfangreiches Werkzeug, um Studenten an Universitäten beim Lernen zu unterstützen. Eine der großen Funktionen von ILIAS ist das Test Modul. Mit Hilfe dieses Moduls können Dozenten Test erstellen, die ihre Studenten beispielsweise lösen müssen, um wichtige Dokumente freizuschalten. Natürlich bietet das Modul auch die Möglichkeit richtige Prüfungen zu erstellen, die dann online stattfinden.

Während das Test Modul von ILIAS nativ bereits eine große Anzahl an verschiedenen Fragetypen unterstützt - wie zum Beispiel Single Choice, Multiple Choice und Sortierungsfragen -, müssen komplexere Fragetypen nachträglich mit Hilfe von Plugins hinzugefügt werden.

In diesem Projekt - genannt \eigenname{\textbf{Q}uestion \textbf{P}lugin for \textbf{I}lias - \textbf{SQL} \textbf{(QPI-SQL)}} - wurde ein ILIAS Plugin - genannt \eigenname{assSQLQuestion} - für diesen Zweck entwickelt. Das Plugin ermöglicht es dabei Dozenten interaktive Fragen zum Themenbereich SQL zu stellen. Studenten können dabei ihre Antworten vor der Abgabe ausführen. Die Abgaben werden dabei außerdem automatisch - durch das Anwenden verschiedener Metriken - bewertet. 
%
% Leere Seite für doppelseitiges plain Layout einfügen
\cleardoubleplainpage
%
% Inhaltsverzeichnis
\tableofcontents
%
% Leere Seite für doppelseitiges plain Layout einfügen
\cleardoubleplainpage
%
% Abbildungsverzeichnis
\listoffigures
%
% Leere Seite für doppelseitiges plain Layout einfügen
\cleardoubleplainpage
%
% Tabellenverzeichnis
\listoftables
%
% Leere Seite für doppelseitiges plain Layout einfügen
\cleardoubleplainpage
%
% Listingverzeichnis
\lstlistoflistings
%
% Leere Seite für doppelseitiges plain Layout einfügen
\cleardoubleplainpage
%
% Seitenlayout mit Seitennummern, Header und Footer
\pagestyle{headings}
%
% Arabische Seitenzahlen
\pagenumbering{arabic}
%
% Seitenzähler auf 1 setzen, das benötigt das Paket hyperref
\setcounter{page}{1}
